\IfFileExists{stacks-project.cls}{%
\documentclass{stacks-project}
}{%
\documentclass{amsart}
}

% For dealing with references we use the comment environment
\usepackage{verbatim}
\newenvironment{reference}{\comment}{\endcomment}
%\newenvironment{reference}{}{}
\newenvironment{slogan}{\comment}{\endcomment}
\newenvironment{history}{\comment}{\endcomment}

% For commutative diagrams we use Xy-pic
\usepackage[all]{xy}

% We use 2cell for 2-commutative diagrams.
\xyoption{2cell}
\UseAllTwocells

% We use multicol for the list of chapters between chapters
\usepackage{multicol}

% This is generally recommended for better output
\usepackage{lmodern}
\usepackage[T1]{fontenc}

% For cross-file-references
\usepackage{xr-hyper}

% Package for hypertext links:
\usepackage{hyperref}

% For any local file, say "hello.tex" you want to link to please
% use \externaldocument[hello-]{hello}
\externaldocument[introduction-]{introduction}
\externaldocument[conventions-]{conventions}
\externaldocument[sets-]{sets}
\externaldocument[categories-]{categories}
\externaldocument[topology-]{topology}
\externaldocument[sheaves-]{sheaves}
\externaldocument[sites-]{sites}
\externaldocument[stacks-]{stacks}
\externaldocument[fields-]{fields}
\externaldocument[algebra-]{algebra}
\externaldocument[brauer-]{brauer}
\externaldocument[homology-]{homology}
\externaldocument[derived-]{derived}
\externaldocument[simplicial-]{simplicial}
\externaldocument[more-algebra-]{more-algebra}
\externaldocument[smoothing-]{smoothing}
\externaldocument[modules-]{modules}
\externaldocument[sites-modules-]{sites-modules}
\externaldocument[injectives-]{injectives}
\externaldocument[cohomology-]{cohomology}
\externaldocument[sites-cohomology-]{sites-cohomology}
\externaldocument[dga-]{dga}
\externaldocument[dpa-]{dpa}
\externaldocument[sdga-]{sdga}
\externaldocument[hypercovering-]{hypercovering}
\externaldocument[schemes-]{schemes}
\externaldocument[constructions-]{constructions}
\externaldocument[properties-]{properties}
\externaldocument[morphisms-]{morphisms}
\externaldocument[coherent-]{coherent}
\externaldocument[divisors-]{divisors}
\externaldocument[limits-]{limits}
\externaldocument[varieties-]{varieties}
\externaldocument[topologies-]{topologies}
\externaldocument[descent-]{descent}
\externaldocument[perfect-]{perfect}
\externaldocument[more-morphisms-]{more-morphisms}
\externaldocument[flat-]{flat}
\externaldocument[groupoids-]{groupoids}
\externaldocument[more-groupoids-]{more-groupoids}
\externaldocument[etale-]{etale}
\externaldocument[chow-]{chow}
\externaldocument[intersection-]{intersection}
\externaldocument[pic-]{pic}
\externaldocument[weil-]{weil}
\externaldocument[adequate-]{adequate}
\externaldocument[dualizing-]{dualizing}
\externaldocument[duality-]{duality}
\externaldocument[discriminant-]{discriminant}
\externaldocument[derham-]{derham}
\externaldocument[local-cohomology-]{local-cohomology}
\externaldocument[algebraization-]{algebraization}
\externaldocument[curves-]{curves}
\externaldocument[resolve-]{resolve}
\externaldocument[models-]{models}
\externaldocument[functors-]{functors}
\externaldocument[equiv-]{equiv}
\externaldocument[pione-]{pione}
\externaldocument[etale-cohomology-]{etale-cohomology}
\externaldocument[proetale-]{proetale}
\externaldocument[relative-cycles-]{relative-cycles}
\externaldocument[more-etale-]{more-etale}
\externaldocument[trace-]{trace}
\externaldocument[crystalline-]{crystalline}
\externaldocument[spaces-]{spaces}
\externaldocument[spaces-properties-]{spaces-properties}
\externaldocument[spaces-morphisms-]{spaces-morphisms}
\externaldocument[decent-spaces-]{decent-spaces}
\externaldocument[spaces-cohomology-]{spaces-cohomology}
\externaldocument[spaces-limits-]{spaces-limits}
\externaldocument[spaces-divisors-]{spaces-divisors}
\externaldocument[spaces-over-fields-]{spaces-over-fields}
\externaldocument[spaces-topologies-]{spaces-topologies}
\externaldocument[spaces-descent-]{spaces-descent}
\externaldocument[spaces-perfect-]{spaces-perfect}
\externaldocument[spaces-more-morphisms-]{spaces-more-morphisms}
\externaldocument[spaces-flat-]{spaces-flat}
\externaldocument[spaces-groupoids-]{spaces-groupoids}
\externaldocument[spaces-more-groupoids-]{spaces-more-groupoids}
\externaldocument[bootstrap-]{bootstrap}
\externaldocument[spaces-pushouts-]{spaces-pushouts}
\externaldocument[spaces-chow-]{spaces-chow}
\externaldocument[groupoids-quotients-]{groupoids-quotients}
\externaldocument[spaces-more-cohomology-]{spaces-more-cohomology}
\externaldocument[spaces-simplicial-]{spaces-simplicial}
\externaldocument[spaces-duality-]{spaces-duality}
\externaldocument[formal-spaces-]{formal-spaces}
\externaldocument[restricted-]{restricted}
\externaldocument[spaces-resolve-]{spaces-resolve}
\externaldocument[formal-defos-]{formal-defos}
\externaldocument[defos-]{defos}
\externaldocument[cotangent-]{cotangent}
\externaldocument[examples-defos-]{examples-defos}
\externaldocument[algebraic-]{algebraic}
\externaldocument[examples-stacks-]{examples-stacks}
\externaldocument[stacks-sheaves-]{stacks-sheaves}
\externaldocument[criteria-]{criteria}
\externaldocument[artin-]{artin}
\externaldocument[quot-]{quot}
\externaldocument[stacks-properties-]{stacks-properties}
\externaldocument[stacks-morphisms-]{stacks-morphisms}
\externaldocument[stacks-limits-]{stacks-limits}
\externaldocument[stacks-cohomology-]{stacks-cohomology}
\externaldocument[stacks-perfect-]{stacks-perfect}
\externaldocument[stacks-introduction-]{stacks-introduction}
\externaldocument[stacks-more-morphisms-]{stacks-more-morphisms}
\externaldocument[stacks-geometry-]{stacks-geometry}
\externaldocument[moduli-]{moduli}
\externaldocument[moduli-curves-]{moduli-curves}
\externaldocument[examples-]{examples}
\externaldocument[exercises-]{exercises}
\externaldocument[guide-]{guide}
\externaldocument[desirables-]{desirables}
\externaldocument[coding-]{coding}
\externaldocument[obsolete-]{obsolete}
\externaldocument[fdl-]{fdl}
\externaldocument[index-]{index}

% Theorem environments.
%
\theoremstyle{plain}
\newtheorem{theorem}[subsection]{Theorem}
\newtheorem{proposition}[subsection]{Proposition}
\newtheorem{lemma}[subsection]{Lemma}

\theoremstyle{definition}
\newtheorem{definition}[subsection]{Definition}
\newtheorem{example}[subsection]{Example}
\newtheorem{exercise}[subsection]{Exercise}
\newtheorem{situation}[subsection]{Situation}

\theoremstyle{remark}
\newtheorem{remark}[subsection]{Remark}
\newtheorem{remarks}[subsection]{Remarks}

\numberwithin{equation}{subsection}

% Macros
%
\def\lim{\mathop{\mathrm{lim}}\nolimits}
\def\colim{\mathop{\mathrm{colim}}\nolimits}
\def\Spec{\mathop{\mathrm{Spec}}}
\def\Hom{\mathop{\mathrm{Hom}}\nolimits}
\def\Ext{\mathop{\mathrm{Ext}}\nolimits}
\def\SheafHom{\mathop{\mathcal{H}\!\mathit{om}}\nolimits}
\def\SheafExt{\mathop{\mathcal{E}\!\mathit{xt}}\nolimits}
\def\Sch{\mathit{Sch}}
\def\Mor{\mathop{\mathrm{Mor}}\nolimits}
\def\Ob{\mathop{\mathrm{Ob}}\nolimits}
\def\Sh{\mathop{\mathit{Sh}}\nolimits}
\def\NL{\mathop{N\!L}\nolimits}
\def\CH{\mathop{\mathrm{CH}}\nolimits}
\def\proetale{{pro\text{-}\acute{e}tale}}
\def\etale{{\acute{e}tale}}
\def\QCoh{\mathit{QCoh}}
\def\Ker{\mathop{\mathrm{Ker}}}
\def\Im{\mathop{\mathrm{Im}}}
\def\Coker{\mathop{\mathrm{Coker}}}
\def\Coim{\mathop{\mathrm{Coim}}}

% Boxtimes
%
\DeclareMathSymbol{\boxtimes}{\mathbin}{AMSa}{"02}

%
% Macros for moduli stacks/spaces
%
\def\QCohstack{\mathcal{QC}\!\mathit{oh}}
\def\Cohstack{\mathcal{C}\!\mathit{oh}}
\def\Spacesstack{\mathcal{S}\!\mathit{paces}}
\def\Quotfunctor{\mathrm{Quot}}
\def\Hilbfunctor{\mathrm{Hilb}}
\def\Curvesstack{\mathcal{C}\!\mathit{urves}}
\def\Polarizedstack{\mathcal{P}\!\mathit{olarized}}
\def\Complexesstack{\mathcal{C}\!\mathit{omplexes}}
% \Pic is the operator that assigns to X its picard group, usage \Pic(X)
% \Picardstack_{X/B} denotes the Picard stack of X over B
% \Picardfunctor_{X/B} denotes the Picard functor of X over B
\def\Pic{\mathop{\mathrm{Pic}}\nolimits}
\def\Picardstack{\mathcal{P}\!\mathit{ic}}
\def\Picardfunctor{\mathrm{Pic}}
\def\Deformationcategory{\mathcal{D}\!\mathit{ef}}
\usepackage{cleveref}
\crefname{originalequation}{original equation}{original equations}
\Crefname{originalequation}{Original equation}{Original equations}
\makeatletter
\newlabel{equation-global-sections-module}{{03GK}{source}{}{sourcepointer-1}{}}
\newlabel{equation-global-sections-module@cref}{{[originalequation][0][]03GK}{[1][0][]source}{}{}{}}
\makeatother
\title{Every quasi-coherent sheaf on projective space is the tilde of its section module}
\author{The Stacks Project Authors}
\date{stacks-project:constructions.tex:local-snapshot-20261012; GNU FDL 1.2 or later}
\begin{document}
\maketitle
\noindent\textit{Editorial note.} 'A small argument, which is omitted here, shows...' covers well-definedness of c across the standard affine cover; the key overlap identity (T\_i/T\_j)\textasciicircum{}d m\_i = m\_j is displayed and argued in-span.\par\medskip
\noindent\textit{Editorial note.} Steps the author leaves to the reader are flagged here for the reviewer: source line(s) 3310 carry an "omitted"/"follows easily"/"immediate from the definitions" phrase. Those sentences are retained verbatim; the displayed argument around them stays inside the counted span.\par\medskip
\noindent\textit{Allowed standard prerequisites (editorial).}\begin{itemize}
\item Quasi-coherent sheaves on affines. schemes-lemma-compare-constructions (cross-chapter, named). Writes the restriction of F to each standard affine as the tilde of a module N\_i.
\end{itemize}\medskip
\ifdefined\hypertarget\hypertarget{sourcepointer-1}{}\fi
\noindent\textit{Source reference.} Original equation 03GK: sources/stacks-constructions/pinned-original-source.tex, source lines 1955--1955. Navigation pointer only: the counted proof cites this statement by \textbackslash{}ref as a standing input; the deduction the counted proof has to make is carried out inside the counted span, and no part of that deduction is delegated to this pointer.\par\medskip
\begin{lemma}
\label{lemma-quasi-coherent-projective-space}
Let $R$ be a ring.
Let $\mathcal{F}$ be a quasi-coherent sheaf on $\mathbf{P}^n_R$.
For $d \geq 0$ set
$$
M_d
=
\Gamma(\mathbf{P}^n_R,
\mathcal{F} \otimes_{\mathcal{O}_{\mathbf{P}^n_R}}
\mathcal{O}_{\mathbf{P}^n_R}(d))
=
\Gamma(\mathbf{P}^n_R, \mathcal{F}(d))
$$
Then $M = \bigoplus_{d \geq 0} M_d$ is a graded $R[T_0, \ldots, R_n]$-module
and there is a canonical isomorphism $\mathcal{F} = \widetilde{M}$.
\end{lemma}

\begin{proof}
The multiplication maps
$$
R[T_0, \ldots, R_n]_e \times M_d \longrightarrow M_{d + e}
$$
come from the natural isomorphisms
$$
\mathcal{O}_{\mathbf{P}^n_R}(e)
\otimes_{\mathcal{O}_{\mathbf{P}^n_R}}
\mathcal{F}(d)
\longrightarrow
\mathcal{F}(e + d)
$$
see Equation (\ref{equation-global-sections-module}). Let us construct the
map $c : \widetilde{M} \to \mathcal{F}$. On each of the standard affines
$U_i = D_{+}(T_i)$ we see that $\Gamma(U_i, \widetilde{M}) = (M[1/T_i])_0$
where the subscript ${}_0$ means degree $0$ part. An element of this
can be written as $m/T_i^d$ with $m \in M_d$. Since $T_i$ is a generator
of $\mathcal{O}(1)$ over $U_i$ we can always write
$m|_{U_i} = m_i \otimes T_i^d$ where $m_i \in \Gamma(U_i, \mathcal{F})$
is a unique section. Thus a natural guess is $c(m/T_i^d) = m_i$.
A small argument, which is omitted here, shows that this gives a
well defined map $c : \widetilde{M} \to \mathcal{F}$ if we can
show that
$$
(T_i/T_j)^d m_i|_{U_i \cap U_j} = m_j|_{U_i \cap U_j}
$$
in $M[1/T_iT_j]$.
But this is clear since on the overlap the generators $T_i$ and
$T_j$ of $\mathcal{O}(1)$ differ by the invertible function $T_i/T_j$.

\medskip\noindent
Injectivity of $c$. We may check for injectivity over the affine opens
$U_i$. Let $i \in \{0, \ldots, n\}$
and let $s$ be an element $s = m/T_i^d \in \Gamma(U_i, \widetilde{M})$
such that $c(m/T_i^d) = 0$. By the description of $c$ above this means that
$m_i = 0$, hence $m|_{U_i} = 0$. Hence $T_i^em = 0$ in $M$ for some
$e$. Hence $s = m/T_i^d = T_i^e/T_i^{e + d} = 0$ as desired.

\medskip\noindent
Surjectivity of $c$. We may check for surjectivity over the affine opens
$U_i$. By renumbering it suffices to check it over $U_0$.
Let $s \in \mathcal{F}(U_0)$.
Let us write $\mathcal{F}|_{U_i} = \widetilde{N_i}$ for some
$R[T_0/T_i, \ldots, T_0/T_i]$-module $N_i$, which is possible because
$\mathcal{F}$ is quasi-coherent. So $s$ corresponds to an element
$x \in N_0$. Then we have that
$$
(N_i)_{T_j/T_i} \cong (N_j)_{T_i/T_j}
$$
(where the subscripts mean ``principal localization at'')
as modules over the ring
$$
R\left[
\frac{T_0}{T_i}, \ldots, \frac{T_n}{T_i},
\frac{T_0}{T_j}, \ldots, \frac{T_n}{T_j}
\right].
$$
This means that for some large integer $d$ there exist elements
$s_i \in N_i$, $i = 1, \ldots, n$ such that
$$
s = (T_i/T_0)^d s_i
$$
on $U_0 \cap U_i$. Next, we look at the difference
$$
t_{ij} = s_i - (T_j/T_i)^d s_j
$$
on $U_i \cap U_j$, $0 < i < j$. By our choice of $s_i$ we know that
$t_{ij}|_{U_0 \cap U_i \cap U_j} = 0$. Hence there exists a large integer
$e$ such that $(T_0/T_i)^et_{ij} = 0$. Set $s_i' = (T_0/T_i)^es_i$,
and $s_0' = s$. Then we will have
$$
s_a' = (T_b/T_a)^{e + d} s_b'
$$
on $U_a \cap U_b$ for all $a, b$. This is exactly the condition that the
elements $s'_a$ glue to a global section
$m \in \Gamma(\mathbf{P}^n_R, \mathcal{F}(e + d))$.
And moreover $c(m/T_0^{e + d}) = s$ by construction. Hence $c$ is
surjective and we win.
\end{proof}

\end{document}
